\section{Procedimentos metodológicos}

O desenho proposto é documental, descritivo, analítico-comparativo e longitudinal, com abordagem mista. A pesquisa preservará uma camada original, com categorias, definições e valores publicados, e uma camada analítica destinada à padronização estrutural. O Sistema de Informações sobre Mortalidade e os registros de segurança pública serão diferenciados dos produtos Anuário Brasileiro de Segurança Pública e Atlas da Violência. A Figura~\ref{fig:arenas-documentais} organiza as arenas de observação previstas no projeto. As relações representadas indicam possibilidades de rastreamento a investigar, sem antecipar uma cadeia causal ou a confirmação de vínculos.

A seleção jornalística abrangerá G1, Folha de S.Paulo e O Estado de S.Paulo, em amostragem intencional associada a eventos de divulgação. A janela principal será do dia da divulgação aos sete dias seguintes, com análises de sensibilidade previstas no projeto. Na arena estatal, serão examinados instrumentos do Legislativo e do Executivo federal, distinguindo citação explícita, correspondência documental e mera proximidade temporal. A comparação entre arenas dialoga com \citeonline{pilar2020atencao}, sem transferir seus resultados norte-americanos para o caso brasileiro. Não se coletarão depoimentos de percepção pública neste desenho.

\begin{figure}[htb]
\caption{Arenas previstas para o rastreamento documental dos indicadores}
\label{fig:arenas-documentais}
\centering
\fbox{\parbox{0.88\textwidth}{\centering
Registros administrativos e construção dos indicadores\\[0.5em]
$\downarrow$\\[0.3em]
Comunicação institucional e seleção jornalística\\[0.5em]
$\downarrow$\\[0.3em]
Mobilização documental no Legislativo e Executivo federal}}
\fonte{Elaboração do autor para o exercício (2026), com base no projeto. As setas não representam efeitos causais comprovados.}
\end{figure}
